\subsection{关于以密度定义的测度的积分}

在本小节的陈述中，g 表示一个正数值函数，它处处有定义且局部 \(\mu\) -可积；\(\theta\) 表示一个复测度，而 u 表示一个局部 \(\theta\) -可积的复函数。

No.~2 中的注记表明，§4 的结果适用于测度 \(\nu = g \cdot \mu = \int g(t)\varepsilon_t d\mu(t)\) （尽管测度 \(g(t)\varepsilon_t\) 仅当 \(g(t) \neq +\infty\) 时才有定义）。这样我们得到以下的陈述：

命题 3. --- 对每个数值函数 \(f \geqslant 0\) （定义在 T 上），

\[
\int^ {\bullet} f d \nu = \int^ {\bullet} (f g) d \mu .\tag{5}
\]

这由 §4, No.~2 的定理 1 得出。

推论 1. --- 为使一个在 Banach 空间中或在 \(\overline{R}\) 中取值的函数 f 对测度 \(u \cdot \theta\) 是局部可忽略的，必须且只须 uf 对 \(\theta\) 是局部可忽略的。

说 f （相应地，\(uf\) ）对 \(u \cdot \theta\) （相应地，对 \(\theta\) ）局部可忽略，等价于说 \(|f|\) （相应地，\(|u| \cdot |f|\) ）对 \(|u \cdot \theta|\) （相应地，对 \(|\theta|\) ）局部可忽略。于是，依据 No.~2 的命题 2，我们被化归到 \(f, u, \theta\) 均为正的情形；这时陈述由命题 3 立即得出。

推论 2. --- 设 \(u_1\) 和 \(u_2\) 是两个局部 \(\theta\) -可积的复函数。为使 \(u_1 \cdot \theta = u_2 \cdot \theta\) ，必须且只须 \(u_1\) 和 \(u_2\) 局部几乎处处相等。

我们立即被化归到证明 \(u \cdot \theta = 0\) 蕴涵 \(u(t) = 0\) 局部几乎处处成立；而 \(u \cdot \theta = 0\) 意味着函数 1 对测度 \(u \cdot \theta\) 局部可忽略。然后应用推论 1 即可。

推论 3. --- 设 u 是一个复函数，它对正测度 \(\mu\) 局部可积。为使 \(u \cdot \mu\) 是一个正测度，必须且只须 \(u(t) \geqslant 0\) 局部几乎处处成立。

因为，\(u \cdot \mu\) 是正的当且仅当 \(u \cdot \mu = |u \cdot \mu| = |u| \cdot \mu\) （命题 2），而这等价于 \(u = |u|\) 局部几乎处处成立（推论 2）。

命题 4. --- 为使一个映射 f （T 到拓扑空间 G 中的）对测度 \(u \cdot \theta\) 可测，必须且只须 f 到 \(\theta\) -可测集 S（由 t 中使 \(u(t) \neq 0\) 者组成）上的限制是 \(\theta\) -可测的。

当 u 和 \(\theta\) 为正时，这由 §4, No.~3 的命题 3 立即得出。然后，借助命题 2，结果推广到 u 和 \(\theta\) 为复的情形。

推论. --- 设 f 是定义在 T 上的一个函数，取值于 Banach 空间 F 或 \(\overline{R}\) 。为使 f 是 \((u\cdot\theta)\) -可测的，必须且只须 uf 是 \(\theta\) -可测的。

因为，\(u\mathbf{f}\) 是 \((uf)|S\) 到 \(T\) 上的用 0 所作的延拓。

定理 1. --- 设 f 是定义在 T 上的一个函数，取值于 Banach 空间 F 或 \(\overline{R}\) 。为使 f 对测度 \(\eta = u\cdot \theta\) 本质可积，必须且只须 uf 本质 \(\theta\) -可积，此时

\[
\int \mathbf {f} d \eta = \int (u \mathbf {f}) d \theta .\tag{6}
\]

再设 \(u\) 连续且 \(u(t) \neq 0\) 对所有 \(t \in \mathbb{T}\) 成立；则 \(\mathbf{f}\) 对测度 \(\eta\) 可积当且仅当 \(u\mathbf{f}\) 对 \(\theta\) 可积。

u 和 \(\theta\) 为正的情形由 §4, No.~4 的定理 2 立即得出。于是陈述的第一个和最后一个断论随之立即成立，因为 f 关于 \(\eta = u \cdot \theta\) 本质可积（相应地，可积）当且仅当它关于 \(|\eta| = |u| \cdot |\theta|\) 本质可积（相应地，可积）。最后，设 f 对 \(\eta\) （从而对 \(|\eta|\) ）本质可积；我们利用分解式 (2)：f 对每个测度 \(\eta_{ij} = g_i \cdot \mu_j\) （i = 1, 2, 3, 4, j = 1, 2, 3, 4）本质可积，因为这些测度都 \(\leqslant |\eta|\) 。我们有

\[
\int \mathbf {f} d \eta_ {i j} = \int g _ {i} \mathbf {f} d \mu_ {j}.
\]

公式 (6) 由此立即得出。

推论. --- 为使测度 \(u \cdot \theta\) 有界，必须且只须 u 本质 \(\theta\) -可积。

例子. --- 设 A 是 T 的一个子集；为使 \(\varphi_{A}\) 局部 \(\mu\) -可积，必须且只须 A 是 \(\mu\) -可测的。假设该条件成立，置 \(\nu = \varphi_{A} \cdot \mu\) ；则对每个数值函数 \(f \geqslant 0\) （定义在 T 上），我们有

\[
\int^ {\bullet} f d \nu = \int^ {\bullet} f \varphi_ {\mathrm{A}} d \mu ,
\]

这个量也记作 \(\int_{\mathrm{A}}^{\bullet}f d\mu\) （参见 Ch. IV, §5, No.~6）。为使一个映射 \(g\) （T 到拓扑空间 G 中的）是 \(\nu\) -可测的，必须且只须 \(g\) 在 A 上的限制是 \(\mu\) -可测的。为使一个映射 \(\mathbf{f}\) （T 到 Banach 空间 F 或 \(\overline{\mathbf{R}}\) 中的）本质 \(\nu\) -可积，必须且只须 \(\mathbf{f}\varphi_{\mathrm{A}}\) 本质 \(\mu\) -可积，此时

\[
\int \mathbf {f} d \nu = \int \mathbf {f} \varphi_ {\mathrm{A}} d \mu ,
\]

这个表达式也记作 \(\int_{A} f d\mu\) 。注意，若 T 到 G （相应地，F、\(\overline{R}\) ）中的两个映射在 A 上一致，则其中一个是 \(\nu\) -可测的（相应地，本质 \(\nu\) -可积的）必须且只须另一个也如此。现在若 g 是 T 的任意子集 \(B \supset A\) 到 G 中的一个映射，则称 g 在 A 上是 \(\mu\) -可测的，如果 g 在 A 上的限制到 T 上的任意延拓是 \(\nu\) -可测的；这等价于说 g 在 A 上的限制是 \(\mu\) -可测的。B 到 Banach 空间 F 或 \(\overline{R}\) 中的一个映射 f 称为在 A 上本质 \(\mu\) -可积的，如果 f 在 A 上的限制到 T 上的某个延拓 \(\overline{f}\) 是本质 \(\nu\) -可积的；此时置

\[
\int_ {\mathrm{A}} \mathbf {f} d \mu = \int_ {\mathrm{A}} \overline {{\mathbf {f}}} d \mu = \int \overline {{\mathbf {f}}} \varphi_ {\mathrm{A}} d \mu ,
\]

并称 \(\int_{\mathrm{A}} \mathbf{f} d\mu\) 为 \(\mathbf{f}\) 在 \(\mathrm{A}\) 上的积分（或延拓到 \(\mathrm{A}\) 的积分）。若 \(f\) 是一个数值函数 \(\geqslant 0\) （定义在 \(\mathrm{B} \supset \mathrm{A}\) 上），则类似地定义 \(\int_{\mathrm{A}}^{*} f d\mu\) 和 \(\int_{\mathrm{A}}^{\bullet} f d\mu\) 。最后，数值函数 \(g\) （定义在 \(\mathrm{B} \supset \mathrm{A}\) 上）称为局部 \(\mu\) -可积于 \(\mathrm{A}\) 的，如果 \(\overline{g}\) （到 \(\mathrm{T}\) 上的一个延拓，由 \(g\) 在 \(\mathrm{A}\) 上的限制扩充而成）是局部 \(\nu\) -可积的：这等价于说，对每个紧子集 \(\mathrm{K}\) （\(\mathrm{T}\) 中的），\(\overline{g}\varphi_{\mathrm{K} \cap \mathrm{A}}\) 是 \(\mu\) -可积的。

